\documentclass[11pt]{article}

\usepackage[margin=1in]{geometry}
\usepackage[T1]{fontenc}
\usepackage[utf8]{inputenc}
\usepackage{amsmath,amssymb}
\usepackage{hyperref}
\usepackage{listings}
\usepackage{xcolor}
\usepackage{longtable}
\usepackage{booktabs}
\usepackage{textcomp}

\hypersetup{colorlinks=true,linkcolor=blue,urlcolor=blue,citecolor=blue}

\definecolor{ctcodebackground}{RGB}{248,248,248}
\definecolor{ctcodekeyword}{RGB}{0,92,175}
\definecolor{ctcodestring}{RGB}{163,21,21}
\definecolor{ctcodecomment}{RGB}{0,128,0}

\lstdefinestyle{ctpython}{
  language=Python,
  basicstyle=\ttfamily\small,
  keywordstyle=\color{ctcodekeyword}\bfseries,
  stringstyle=\color{ctcodestring},
  commentstyle=\color{ctcodecomment}\itshape,
  backgroundcolor=\color{ctcodebackground},
  breaklines=true,
  columns=fullflexible,
  frame=single,
  framerule=0.4pt,
  keepspaces=true,
  showspaces=false,
  showstringspaces=false,
  showtabs=false,
  numbers=none,
  tabsize=4,
  upquote=true
}

\lstdefinestyle{ctterminal}{
  basicstyle=\ttfamily\small,
  backgroundcolor=\color{ctcodebackground},
  breaklines=true,
  columns=fullflexible,
  frame=single,
  framerule=0.4pt,
  keepspaces=true,
  showspaces=false,
  showstringspaces=false,
  showtabs=false,
  numbers=none,
  tabsize=4,
  upquote=true
}

\lstset{style=ctpython}

\title{\texttt{xstar-atomic} User Guide}
\author{XSTAR ATDB direct decoder, source-code-aligned rate audits, and prototype solvers}
\date{Version 0.3.129}

\begin{document}
\maketitle
\tableofcontents
\newpage

\section{Overview}

\texttt{xstar-atomic} reads XSTAR's packed \texttt{atdb.fits} atomic database, evaluates selected XSTAR rate formulae, builds prototype level-population products, and provides source-code-first validation tools for comparing against same-run XSTAR outputs.

The package is not intended to replace XSTAR. Its purpose is to make the atomic-data and local-rate pieces auditable from Python, with explicit provenance for the XSTAR record, source-code branch, local plasma state, radiation field, escape treatment, and population-matrix term.

\begin{lstlisting}[style=ctpython]
import xstar_atomic as xa
from xstar_atomic import XSTARAtomic
\end{lstlisting}

\section{Installation and data setup}

\subsection{Install the package}

From the source tree:

\begin{lstlisting}[style=ctterminal]
python -m pip install -e .
\end{lstlisting}

For development and tests:

\begin{lstlisting}[style=ctterminal]
python -m pip install -e .[dev]
PYTHONPATH=src pytest -q
\end{lstlisting}

Optional extras:

\begin{lstlisting}[style=ctterminal]
python -m pip install -e .[docs]
python -m pip install -e .[hdf5]
python -m pip install -e .[sparse]
\end{lstlisting}

\subsection{Configure \texttt{atdb.fits}}

The main database is XSTAR's packed atomic database, usually located at:

\begin{lstlisting}[style=ctterminal]
xstar/data/atdb.fits
\end{lstlisting}

You can provide it explicitly:

\begin{lstlisting}[style=ctpython]
db = XSTARAtomic("/path/to/xstar/data/atdb.fits")
\end{lstlisting}

or configure a data path:

\begin{lstlisting}[style=ctpython]
import xstar_atomic as xa

xa.set_data_path("/path/to/xstar/data/atdb.fits")
print(xa.get_data_path())
\end{lstlisting}

\section{Quick start}

\subsection{Open the database and inspect the hierarchy}

\begin{lstlisting}[style=ctpython]
from xstar_atomic import XSTARAtomic

db = XSTARAtomic("/path/to/xstar/data/atdb.fits")
print(db.summary())
\end{lstlisting}

\subsection{Extract levels and lines}

\begin{lstlisting}[style=ctpython]
levels = db.levels("O VII")
lines = db.lines("O VII", wavelength=(21.4, 22.2), slim=True)
\end{lstlisting}

\subsection{Evaluate collision and emissivity products}

\begin{lstlisting}[style=ctpython]
coll = db.collisions("O VIII", temperatures=[1e6, 3e6], wavelength=(18.8, 19.1))
emiss = db.emissivity("O VIII", temperatures=[1e6, 3e6], wavelength=(18.8, 19.1))
\end{lstlisting}

\subsection{Build an audit context and evaluate a type-50 rate}

The v0.3.128 API infrastructure adds explicit context objects. These are required because XSTAR rates can depend on the local state and radiation/escape terms, not only on the atomic record.

\begin{lstlisting}[style=ctpython]
from xstar_atomic import LocalPlasmaState, EscapeContext, evaluate_type50_bound_bound

state = LocalPlasmaState(
    temperature_K=7.66552e4,
    electron_density_cm3=1.20466e8,
    log_xi=1.5,
)

escape = EscapeContext(
    cfrac=0.0,
    ptmp1=0.2,
    ptmp2=0.3,
    flinabs_ptmp1=0.8,
)

rate = evaluate_type50_bound_bound(
    aij_s_inv=1.0e12,
    oscillator_strength=0.7,
    wavelength_A=21.602,
    vtherm_cm_s=1.0e7,
    bremsa_nb1=1.0e5,
    plasma_state=state,
    escape_context=escape,
    ion="O VII",
)

print(rate.status)
print(rate.lower_to_upper_photoexcitation_s_inv)
print(rate.upper_to_lower_escaped_decay_s_inv)
print(rate.source_formula)
\end{lstlisting}

In v0.3.128 this evaluator is audit-only. It records the XSTAR \texttt{ucalc.f90} type-50 branch and returns a \texttt{RateEvaluation} object; it does not inject photoexcitation into the population solver.


\section{Public API cookbook}

This section summarizes the most useful public API calls. The goal is to make common workflows discoverable without reading the internal example scripts.

\subsection{Open and summarize an XSTAR database}

\begin{lstlisting}[style=ctpython]
import xstar_atomic as xa
from xstar_atomic import XSTARAtomic

xa.set_data_path("/path/to/xstar/data/atdb.fits")
db = XSTARAtomic()             # uses the configured datapath
summary = db.summary()
print(summary)
\end{lstlisting}

\subsection{Query levels, lines, and wavelength windows}

\begin{lstlisting}[style=ctpython]
levels = db.levels("O VII")
triplet_lines = db.lines("O VII", wavelength=(21.4, 22.2), slim=True)
lya = db.lines("O VIII", wavelength=(18.8, 19.1), slim=True)

for line in triplet_lines:
    print(line.get("wavelength_A"), line.get("lower_level"), line.get("upper_level"))
\end{lstlisting}

Use this layer for quick inspection and for building small validation tables. Use lower-level record access only when raw ATDB indices are required.

\subsection{Evaluate collisional and recombination products}

\begin{lstlisting}[style=ctpython]
coll = db.collisions(
    "O VIII",
    temperatures=[1.0e6, 3.0e6, 1.0e7],
    wavelength=(18.8, 19.1),
)

rr = db.recombination(
    "O VII",
    temperatures=[1.0e6],
    electron_densities=[1.0e8],
    source_mode="none",
)

print(len(coll["summary"]), len(coll["evaluated"]))
print(rr["summary"])
\end{lstlisting}

\subsection{Build a local XSTAR-style context}

\begin{lstlisting}[style=ctpython]
from xstar_atomic import LocalPlasmaState, RadiationField, EscapeContext

state = LocalPlasmaState(
    temperature_K=7.66552e4,
    electron_density_cm3=1.20466e8,
    log_xi=1.5,
    ion_fraction=0.255926,
)

rad = RadiationField.from_pairs([
    (0.50, 1.0e4),
    (0.574, 2.5e4),
    (1.00, 1.0e4),
])

escape = EscapeContext(
    cfrac=0.0,
    ptmp1=0.2,
    ptmp2=0.3,
    flinabs_ptmp1=0.8,
)
\end{lstlisting}

These objects are intentionally explicit. XSTAR rate terms can depend on local temperature, electron density, radiation-field binning, covering factor, and escape probabilities, so those values should not be hidden inside scalar helper functions.

\subsection{Evaluate the audit-only type-50 bound-bound branch}

\begin{lstlisting}[style=ctpython]
from xstar_atomic import evaluate_type50_bound_bound

rate = evaluate_type50_bound_bound(
    aij_s_inv=3.0e12,
    oscillator_strength=0.7,
    wavelength_A=21.602,
    vtherm_cm_s=1.0e7,
    bremsa_nb1=2.5e4,
    plasma_state=state,
    radiation_field=rad,
    escape_context=escape,
    ion="O VII",
    lower_level=1,
    upper_level=7,
)

print(rate.lower_to_upper_photoexcitation_s_inv)
print(rate.upper_to_lower_escaped_decay_s_inv)
print(rate.terms)
\end{lstlisting}

In v0.3.128 and v0.3.129 this evaluator remains audit-only. It records the XSTAR \texttt{ucalc.f90} type-50 branch and branch swap but does not modify the level-population solver.

\subsection{Use high-level convenience methods}

\begin{lstlisting}[style=ctpython]
rate = db.type50_rate(
    aij_s_inv=3.0e12,
    oscillator_strength=0.7,
    wavelength_A=21.602,
    vtherm_cm_s=1.0e7,
    bremsa_nb1=2.5e4,
    plasma_state=state,
    escape_context=escape,
    ion="O VII",
)
\end{lstlisting}

The method delegates to the same source-code-aligned evaluator, so it is convenient for notebooks while preserving provenance.

\subsection{Run a source-code audit from Python}

\begin{lstlisting}[style=ctpython]
from xstar_atomic.audit import type50_line_pumping

audit = type50_line_pumping(
    cases_csv="helike_local_state_validation_v03129/helike_local_state_cases.csv",
    solver_root=".",
    xstar_source_root="../xstar",
    out_dir="helike_type50_line_pumping_audit_v03129",
    print_summary=True,
)

print(audit.summary)
\end{lstlisting}

The corresponding command-line wrapper is \texttt{examples/55\_audit\_helike\_type50\_line\_pumping.py}.

\subsection{Export line-emissivity products}

\begin{lstlisting}[style=ctpython]
from xstar_atomic.export import export_superwind_bundle, parse_band_specs

bands = parse_band_specs(["soft:0.5:2.0", "hard:2.0:10.0"])
bundle = export_superwind_bundle(
    "/path/to/atdb.fits",
    ions="O VIII,Ne IX",
    out_dir="atomic_export_example",
    temperatures=[1.0e6, 3.0e6, 1.0e7],
    wavelength=(1.0, 40.0),
    formats=["csv", "hdf5"],
    bands=bands,
)
\end{lstlisting}

This is the current public export workflow for downstream simulation or spectral-model post-processing.

\section{Atomic database access}

\subsection{Packed XSTAR database structure}

The public XSTAR \texttt{atdb.fits} file contains packed arrays rather than one simple table:

\begin{lstlisting}[style=ctterminal]
POINTERS
REALS
INTEGERS
CHARS
\end{lstlisting}

\texttt{xstar-atomic} rebuilds the hierarchy:

\begin{lstlisting}[style=ctterminal]
element -> ion -> levels -> radiative/collisional/photoionization/recombination records
\end{lstlisting}

The low-level reader is \texttt{ATDB}; the higher-level convenience class is \texttt{XSTARAtomic}.

\subsection{Public database functions}

Common functions and methods include:

\begin{lstlisting}[style=ctpython]
xa.find_atdb_file(...)
xa.resolve_atdb_path(...)
xa.set_data_path(...)
xa.get_data_path()

XSTARAtomic(...).summary()
XSTARAtomic(...).levels("O VII")
XSTARAtomic(...).lines("O VII")
XSTARAtomic(...).photoionization("O VII")
XSTARAtomic(...).collisions("O VII")
XSTARAtomic(...).recombination("O VII")
XSTARAtomic(...).emissivity("O VII")
\end{lstlisting}

\subsection{Current decoder status}

\begin{longtable}{p{0.35\linewidth}p{0.18\linewidth}p{0.37\linewidth}}
\toprule
Component & XSTAR data type(s) & Current status \\
\midrule
Packed FITS reader / hierarchy & arrays & validated on released \texttt{atdb.fits} \\
Levels & 6 & implemented \\
Radiative lines & 50 & implemented; v0.3.128 adds an audit-only source-code rate object for line pumping \\
Photoionization grids & 53 & implemented for ordinary OP-style grids \\
Collisions, tabulated upsilon & 56 & implemented \\
Collisions, Bautista hydrogenic/l-mixing branches & 63 & implemented for validated branches \\
Collisions, Burgess--Tully & 51, 98 & implemented for representative branches \\
Recombination totals & 1, 30, 38, 39 & implemented for total rates \\
Charge exchange & 2 & implemented when explicitly requested \\
Full local radiation-field parity & several & in development \\
Production replacement for XSTAR transfer/thermal balance & combined & not a package goal \\
\bottomrule
\end{longtable}

\section{Local context objects}

\subsection{\texttt{LocalPlasmaState}}

\texttt{LocalPlasmaState} stores local gas quantities:

\begin{lstlisting}[style=ctpython]
from xstar_atomic import LocalPlasmaState

state = LocalPlasmaState(
    temperature_K=1.0e6,
    electron_density_cm3=1.0e8,
    log_xi=3.0,
    ion_fraction=0.25,
)
\end{lstlisting}

A state can also be built from a row dictionary, including rows from local-state validation CSV files:

\begin{lstlisting}[style=ctpython]
state = LocalPlasmaState.from_mapping(row)
\end{lstlisting}

\subsection{\texttt{RadiationField}}

\texttt{RadiationField} stores an XSTAR-like local radiation field and energy grid:

\begin{lstlisting}[style=ctpython]
from xstar_atomic import RadiationField

rad = RadiationField.from_pairs([(0.5, 1.0e4), (1.0, 2.0e4)])
print(rad.value_at(0.9))
\end{lstlisting}

The current nearest-bin helper is only an audit convenience. Exact XSTAR \texttt{nbinc} parity is a future development target.

\subsection{\texttt{EscapeContext}}

\texttt{EscapeContext} stores line escape and covering-factor terms:

\begin{lstlisting}[style=ctpython]
from xstar_atomic import EscapeContext

escape = EscapeContext(cfrac=0.0, ptmp1=0.1, ptmp2=0.2, flinabs_ptmp1=0.7)
\end{lstlisting}

For type-50 line pumping, the upward photoexcitation branch is multiplied by \texttt{max(0, 1-cfrac)}.

\section{Source-code-aligned rate evaluators}

\subsection{\texttt{RateEvaluation}}

\texttt{RateEvaluation} is a structured result object. It records process, data type, status, value, source formula, record id, ion, lower and upper levels, terms, warnings, and metadata. This is intentionally more verbose than a simple scalar rate. It supports source-code-first audits and future matrix-term provenance.

\subsection{Type-50 bound-bound radiative branch}

XSTAR's type-50 branch forms two local quantities before a final branch swap:

\begin{lstlisting}[style=ctterminal]
ans1 = aij * (ptmp1 + ptmp2)
sigma = 0.02655 * flin * elin * 1.d-8 / vtherm
ans2 = sigma * bremsa(nb1) * vtherm / 3.e10 * flinabs(ptmp1)
ans2 = ans2 * max(0., 1.d0-cfrac)
\end{lstlisting}

After the swap:

\begin{lstlisting}[style=ctterminal]
ans1 -> lower-to-upper photoexcitation
ans2 -> upper-to-lower escaped decay
\end{lstlisting}

The v0.3.128 function is:

\begin{lstlisting}[style=ctpython]
from xstar_atomic import evaluate_type50_bound_bound

rate = evaluate_type50_bound_bound(...)
\end{lstlisting}

This function is audit-only and should be used to inspect the missing line-pumping term before any future solver-injection mode is enabled.

\section{Matrix, solver, and emissivity workflows}

The existing solver remains prototype/diagnostic. Current public methods include:

\begin{lstlisting}[style=ctpython]
db.emissivity(...)
db.type50_rate(...)
db.audit_type50_line_pumping(...)
\end{lstlisting}

Future stable APIs should separate:

\begin{lstlisting}[style=ctterminal]
rate evaluator -> matrix term -> population solution -> emissivity/triplet result
\end{lstlisting}

Each matrix term should carry provenance: XSTAR data type, record id, source-code formula, local context terms, and treatment mode.

\section{XSTAR-output readers and same-run validation}

\texttt{xstar\_atomic.xstar\_outputs} reads selected XSTAR output FITS files, including local abundance/state tables and line lists. The local-state validation examples use these readers to select the same local temperature, electron density, and same-run triplet targets before comparing solver output.

\begin{lstlisting}[style=ctterminal]
PYTHONPATH=src python examples/51_run_helike_local_state_validation.py \
  --xstar-runs-root xstar_runs \
  --target-root xstar_atomic_v0.3.111_results \
  --atdb ../xstar/data/atdb.fits \
  --selection-mode max \
  --target-electron-density 1e8 \
  --nearest-density \
  --out-dir helike_local_state_validation_v03129_ne1e8_allions \
  --print-summary
\end{lstlisting}

\section{Audit workflows}

\subsection{Type-50 line-pumping audit}

In v0.3.128 the previous example-55 logic is available as a Python API:

\begin{lstlisting}[style=ctpython]
from xstar_atomic.audit import type50_line_pumping

audit = type50_line_pumping(
    "helike_local_state_validation_v03129_ne1e8_allions/helike_local_state_cases.csv",
    solver_root=".",
    xstar_source_root="../xstar",
    out_dir="helike_type50_line_pumping_audit_v03129",
    print_summary=True,
)
\end{lstlisting}

The CLI wrapper remains:

\begin{lstlisting}[style=ctterminal]
PYTHONPATH=src python examples/55_audit_helike_type50_line_pumping.py \
  --cases-csv helike_local_state_validation_v03129_ne1e8_allions/helike_local_state_cases.csv \
  --solver-root . \
  --xstar-source-root ../xstar \
  --out-dir helike_type50_line_pumping_audit_v03129 \
  --print-summary
\end{lstlisting}

The audit writes:

\begin{lstlisting}[style=ctterminal]
helike_type50_line_pumping_audit.csv
helike_type50_line_pumping_audit.md
helike_type50_line_pumping_audit.json
\end{lstlisting}

\subsection{Audit policy}

Development policy:

\begin{enumerate}
\item No empirical triplet scale factors as final physics.
\item New physics starts as audit-only.
\item Solver-changing modes must be opt-in until validated.
\item Every solver-changing rate must have source-code and record provenance.
\item Radiation-dependent rates must expose the radiation context.
\item Escape-dependent rates must expose \texttt{cfrac}, \texttt{tau/ptmp}, and escape treatment.
\item Same-run XSTAR outputs are preferred over generic target CSVs.
\end{enumerate}

\section{Examples and source-module migration}

The package contains many examples because the XSTAR-alignment work is still exploratory. As workflows stabilize, reusable logic should move into \texttt{src/xstar\_atomic/}, leaving examples as short wrappers.

\begin{longtable}{p{0.42\linewidth}p{0.48\linewidth}}
\toprule
Example(s) & Candidate source module/API \\
\midrule
\texttt{51\_run\_helike\_local\_state\_validation.py} & \texttt{xstar\_atomic.runs.select\_local\_states(...)} \\
\texttt{52\_summarize\_helike\_local\_state\_comparison.py} & \texttt{xstar\_atomic.validate.summarize\_local\_state\_comparison(...)} \\
\texttt{53\_audit\_helike\_resonance\_deficit.py} & \texttt{xstar\_atomic.audit.resonance\_deficit(...)} \\
\texttt{54\_audit\_helike\_resonance\_population\_flux.py} & \texttt{xstar\_atomic.audit.resonance\_population\_flux(...)} \\
\texttt{55\_audit\_helike\_type50\_line\_pumping.py} & \texttt{xstar\_atomic.audit.type50\_line\_pumping(...)} \\
\texttt{43\_compare\_xstar\_detail\_populations.py} & \texttt{xstar\_atomic.validate.compare\_detail\_populations(...)} \\
\texttt{42\_xstar\_like\_element\_solver\_demo.py} & \texttt{xstar\_atomic.solve.element(...)} \\
\texttt{47\_prepare\_mg\_ca\_xstar\_triplet\_targets.py} & \texttt{xstar\_atomic.xout.convert\_triplet\_targets(...)} \\
\bottomrule
\end{longtable}

See \texttt{examples/README.md} for grouped example usage and \texttt{docs/example\_to\_source\_api\_map.md} for a broader migration plan.

\section{Validation and tests}

Recommended local checks:

\begin{lstlisting}[style=ctterminal]
PYTHONPATH=src python -m compileall -q src examples tests
PYTHONPATH=src pytest -q tests/test_api_infrastructure_v03128.py
PYTHONPATH=src pytest -q \
  tests/test_type50_line_escape.py \
  tests/test_helike_triplet_balance_diagnostics.py \
  tests/test_xstar_detail_population_compare_example.py \
  tests/test_package_metadata.py \
  tests/test_cli_smoke.py
\end{lstlisting}

Some tests are skipped unless \texttt{XSTAR\_ATDB\_FITS} points to a local \texttt{atdb.fits}.

\section{Roadmap}

Near-term:

\begin{lstlisting}[style=ctterminal]
v0.3.128: API infrastructure and audit-only type-50 evaluator.
v0.3.129: examples README, expanded API cookbook, and LaTeX table of contents/documentation polish.
v0.3.130: XSTAR radiation-field reader and exact line-energy bin mapping.
v0.3.131: type-50 matrix-injection preview, still audit-only.
v0.3.132: optional solver mode for xstar-line-escape-and-pumping.
\end{lstlisting}

Longer-term:

\begin{lstlisting}[style=ctterminal]
- stable context/rate/matrix/solve/emissivity APIs
- broader data-type evaluator parity
- stronger same-run XSTAR validation suite
- source-code provenance for every matrix term
- documentation organized by workflow rather than version history
\end{lstlisting}

\end{document}
